\begin{abstract}
Advancements like BitNet~\cite{bitnet} have initiated a transformative era for 1-bit Large Language Models (LLMs). In this paper, we present a 1-bit LLM variant, referred to as \textbf{\text{BitNet b1.58}}, characterized by having each model parameter drawn from the ternary set \{-1, 0, 1\}. Remarkably, this model achieves parity with standard full-precision Transformers (e.g., FP16 or BF16) of equivalent model size and trained with the same number of tokens, not only in terms of perplexity but also on downstream tasks, while offering dramatic improvements in latency, memory footprint, throughput, and energy efficiency. Notably, the 1.58-bit LLM introduces a new scaling paradigm, establishing a blueprint for constructing future LLMs that reconcile high performance with resource efficiency. This innovation also establishes a foundation for new computational paradigms and unlocks opportunities to develop specialized hardware tailored for 1-bit LLM architectures.
\end{abstract}

\section{The Era of 1-bit LLMs}

Artificial intelligence research has recently experienced significant expansion in both the scale and proficiency of Large Language Models (LLMs). While these models excel across an array of natural language understanding and generation benchmarks, their ever-increasing parameter counts pose formidable obstacles to practical deployment, fueling apprehensions regarding their substantial resource demands and environmental footprint due to elevated power consumption.

To alleviate these burdens, post-training quantization has emerged as a practical solution, producing low-precision models targeted for efficient inference~\cite{smoothquant, gptq, quip, quip_sharp}. By quantizing both weights and activations, this method markedly diminishes the memory and computational requirements inherent to LLMs. The community trend has been a progression from 16-bit formats to lower-precision schemes, such as 4-bit implementations~\cite{gptq, awq}. Nevertheless, although post-training quantization is a widely adopted technique within the industry, it is inherently sub-optimal compared to alternative strategies.

Recent advancements in 1-bit model architectures, exemplified by BitNet~\cite{bitnet}, indicate a promising strategy for decreasing the operational costs of large language models (LLMs) without sacrificing effectiveness. Standard LLMs typically utilize 16-bit floating point numbers (such as FP16 or BF16), with matrix multiplication constituting the primary computation workload. Consequently, most computational expense arises from floating-point addition and multiplication. BitNet, by contrast, implements matrix multiplication that relies solely on integer addition, which drastically lowers energy consumption requirements for LLMs. Since power constraints often cap the computing throughput of contemporary chips, reduced energy usage directly contributes to increased computational speed.

Beyond computation, transferring model weights from DRAM to on-chip memory (for example, SRAM) during inference imposes significant overhead. While expanding SRAM can theoretically enhance throughput, such modifications lead to considerable expense, surpassing that of DRAM. 1-bit LLMs, when compared to their full-precision counterparts, are far more economical with respect to both memory capacity and bandwidth. This reduced memory demand allows for quicker, cheaper loading of model weights from DRAM, thereby boosting inference efficiency.

In this paper, we propose a major 1-bit LLM modification: \textbf{\text{BitNet b1.58}}, wherein every parameter adopts a ternary value from the set \{-1, 0, 1\}. By incorporating zero into the original BitNet’s parameter set, we achieve an effective storage size of 1.58 bits per parameter. \text{BitNet b1.58} preserves the benefits seen in the original 1-bit BitNet, notably its computational model that forgoes multiplication in favor of addition for matrix operations—this paradigm is readily optimized for performance and remains highly energy-efficient. In addition, it demonstrates substantial improvements in memory usage, throughput, and latency compared to FP16 LLM benchmarks. Beyond these similarities, \text{BitNet b1.58} confers two distinct benefits. First, the inclusion of 0 explicitly enables feature filtering within model weights, thereby enhancing the expressiveness and effectiveness of the 1-bit LLM. Second, empirical results reveal that, for models of size 3B or larger and under identical configurations (e.g., parameter count, training data), \text{BitNet b1.58} achieves perplexity and downstream task performance on par with full-precision (FP16) baselines.

\section{BitNet b1.58}

\text{BitNet b1.58} builds upon the BitNet framework, a variant of the Transformer in which the standard \emph{nn.Linear} operations are replaced by \emph{BitLinear}. The model is initialized and trained from the beginning using weights quantized to 1.58 bits and activations quantized to 8 bits. Several adaptations distinguish it from the original BitNet, which we outline below.

\paragraph{Quantization Function.} In order to restrict the weights to the discrete set $\{-1, 0, +1\}$, we employ an \emph{absmean} quantization method. This approach involves normalizing the weight matrix by the mean of its absolute values before quantizing each entry to its closest integer within $\{-1, 0, +1\}$:

\begin{equation}
    \widetilde{W} = \mathrm{RoundClip}(\frac{W}{\gamma + \epsilon}, -1, 1),
\end{equation}

\begin{equation}
    \mathrm{RoundClip}(x, a, b) = \max(a, \min(b, \mathrm{round}(x))),
\end{equation}

\begin{equation}
    \gamma = \frac{1}{nm}\sum_{ij} |W_{ij}|.
\end{equation}

The process for activation quantization retains the same structure as utilized in BitNet. However, unlike BitNet, we do not perform rescaling of activations into the interval $[0, Q_b]$ prior to applying non-linear transformations. Instead, activations are uniformly rescaled within $[-Q_b, Q_b]$ for each token, effectively eliminating zero-point quantization. This revision streamlines both software implementation and hardware optimization, and our empirical analysis indicates it yields an insignificant impact on overall model performance.

\paragraph{LLaMA-alike Components.}

LLaMA~\cite{llama, llama2} has emerged as the standard reference architecture among open-source LLM implementations. To align with and facilitate broader adoption within the open-source community, \text{BitNet b1.58} incorporates LLaMA-inspired architectural elements, including RMSNorm~\cite{rmsnorm}, SwiGLU~\cite{swiglu}, rotary embedding~\cite{rope}, and the removal of all bias terms. This architectural compatibility allows \text{BitNet b1.58} to be readily deployed in widely used open-source toolchains, such as Huggingface, vLLM~\cite{vllm}, and llama.cpp\footnote{\url{https://github.com/ggerganov/llama.cpp}}, with minimal adaptation required.

\section{Results}

We benchmarked \text{BitNet b1.58} against our reproduced FP16 \text{LLaMA LLM} variants of multiple sizes. Both model families were pre-trained from scratch using 100 billion tokens sampled from the RedPajama dataset~\cite{redpajama} to ensure experimental consistency. Zero-shot evaluation was conducted on an array of language understanding benchmarks: ARC-Easy~\cite{arc}, ARC-Challenge~\cite{arc}, Hellaswag~\cite{hellaswag}, Winogrande~\cite{winoGrande}, PIQA~\cite{piqa}, OpenbookQA~\cite{openbookqa}, and BoolQ~\cite{boolq}. Additionally, we calculated validation perplexity using the WikiText2~\cite{wikitext2} and C4~\cite{c4} corpora.

The GPU memory consumption and inference latency of both \text{LLaMA LLM} and \text{BitNet b1.58} were compared by deploying them through the FasterTransformer\footnote{\url{https://github.com/NVIDIA/FasterTransformer}} framework, a system specifically tailored to minimize inference latency of LLMs on GPUs. For \text{BitNet b1.58}, the 2-bit Ladder~\cite{ladder} kernel was incorporated. Our evaluation metric for runtime cost was the per-token output latency, reflecting the principal component of inference expenses.

The key findings are summarized in Table~\ref{tab:ppl}, which displays both perplexity and operational resource usage for \text{BitNet b1.58} and \text{LLaMA LLM}. Notably, starting at the 3B parameter scale, \text{BitNet b1.58} matches the perplexity of full-precision \text{LLaMA LLM} while providing 2.71-fold speedup and a 3.55-fold reduction in GPU memory consumption. The 3.9B version of \text{BitNet b1.58} offers 2.4x faster inference, 3.32x lower memory usage, and also surpasses the 3B \text{LLaMA LLM} in task performance.

Comprehensive zero-shot accuracy across downstream tasks is presented in Table~\ref{tab:zero-shot}. Following the \emph{lm-evaluation-harness}\footnote{\url{https://github.com/EleutherAI/lm-evaluation-harness}} protocol, we find that as model size increases, the performance differential between \text{BitNet b1.58} and \text{LLaMA LLM} diminishes. Most significantly, \text{BitNet b1.58} attains parity with its full-precision counterpart at the 3B scale. Consistent with perplexity trends, the 3.9B \text{BitNet b1.58} variant outperforms the 3B \text{LLaMA LLM} while maintaining considerable advantages in both memory and inference latency. These results establish \text{BitNet b1.58} as a Pareto optimal alternative within the latest generation of LLM architectures.

\paragraph{Memory and Latency}

We expanded the model scale to 7B, 13B, and 70B parameters and systematically analyzed the associated computational costs. As shown in Figure~\ref{fig:latency-memory-energy}, both latency and memory usage reveal clear patterns as model size increases, with acceleration benefits becoming more pronounced. Notably, \text{BitNet b1.58} at 70B achieves a speed-up of 4.1$\times$ relative to the \text{LLaMA LLM} reference. This improvement can be attributed to the growing computational expense of the \emph{nn.Linear} operation in larger models. Memory utilization exhibits a similar trajectory, given that the embedding layer, which retains full precision, constitutes a smaller fraction of overall memory in larger architectures. Both the reported latency and memory metrics were measured using a 2-bit kernel, indicating that further gains in cost efficiency are feasible through additional optimizations.

\paragraph{Energy}

We also carried out estimates of arithmetic energy requirements for \text{BitNet b1.58} and \text{LLaMA LLM}, focusing principally on the matrix multiplication stage, which dominates energy consumption in LLM computations. Figure~\ref{fig:energy} provides a breakdown of these energy costs: INT8 additions form the bulk of \text{BitNet b1.58}'s energy expenditure, whereas \text{LLaMA LLM} predominantly consumes energy through both FP16 additions and multiplications. Citing the energy analysis framework in~\cite{energycost, pokebnn}, we observe that \text{BitNet b1.58} reduces matrix multiplication energy consumption by 71.4$\times$ on 7nm technology nodes. Additionally, we quantified end-to-end energy usage for models processing 512 tokens, observing that the energy efficiency advantage of \text{BitNet b1.58} relative to the FP16 \text{LLaMA LLM} continues to grow with model size. This arises from the fact that as model scale increases, the proportion of computation spent on \emph{nn.Linear} grows, whereas other components account for a diminishing share of overall cost in larger models.

\paragraph{Throughput}

We evaluate the throughput of \text{BitNet b1.58} in comparison to the \text{LLaMA LLM} 70B on a setup with two A100 GPUs (each with 80GB), leveraging pipeline parallelism~\cite{gpipe} to accommodate the size of the \text{LLaMA LLM} 70B on this hardware. The batch size is progressively increased until it hits the memory limit of the GPUs, fixing the sequence length at 512. As illustrated in Table~\ref{tab:throughput}, \text{BitNet b1.58} 70B achieves batch sizes that are up to 11 times larger than those supported by \text{LLaMA LLM}, yielding a throughput increase by a factor of 8.9.

\textbf{\text{BitNet b1.58} establishes a new scaling relationship linking model capacity, performance, and inference efficiency}. Drawing on the data from Figure~\ref{fig:latency-memory-energy} and \ref{fig:energy}, we outline the following approximate correspondences in efficiency between 1.58-bit and 16-bit models:

\begin{itemize}
    \item A 13B parameter BitNet b1.58 surpasses a 3B FP16 LLM in latency, memory footprint, and power consumption.
    \item A 30B parameter BitNet b1.58 outperforms a 7B FP16 LLM regarding latency, memory demand, and energy usage.
    \item A 70B parameter BitNet b1.58 exceeds the efficiency—across latency, memory, and energy—of a 13B FP16 LLM.
\end{itemize}

\paragraph{Training with 2T Tokens}

The total count of training tokens is a vital consideration for scaling large language models. To examine how well \text{BitNet b1.58} scales with respect to token volume, we trained a version of the model on 2T tokens following the data selection procedure established for StableLM-3B~\cite{StableLM-3B-4E1T}, currently the leading open 3B-parameter model. Evaluation was conducted on a suite comprising Winogrande~\cite{winoGrande}, PIQA~\cite{piqa}, SciQ~\cite{sciq}, LAMBADA~\cite{lambada}, and ARC-easy~\cite{arc}. Zero-shot accuracy results are compiled in Table~\ref{tab:2t}. For datasets using both accuracy and normalized accuracy metrics, their mean is reported. Accuracy results for StableLM 3b after 2T token training are taken from the corresponding technical documentation. Our analysis reveals that \text{BitNet b1.58} consistently delivers stronger results across all evaluated benchmarks, highlighting the robust generalization achieved by the 1.58-bit models.

\section{Discussion and Future Work}

\textbf{1-bit Mixture-of-Experts (MoE) LLMs}

Mixture-of-Experts (MoE) architectures have established themselves as an efficient solution for scaling LLMs by curbing computational FLOPs. Nevertheless, excessive memory usage and significant inter-chip communication remain major bottlenecks for practical MoE deployment. The advent of 1.58-bit LLMs offers a promising solution to these issues. The sharply reduced memory demands decrease the need for distributed devices to host MoE models, and minimize the communication cost associated with transferring activations. In optimal scenarios, if entire models fit within a single chip, the issue of inter-chip transfer overhead could be completely eliminated.

\textbf{Native Support of Long Sequence in LLMs}

Long sequence processing has emerged as a central requirement in large language models. However, memory consumption from the storage of KV caches poses a significant challenge for inference over extended contexts. The \text{BitNet b1.58} model marks considerable progress by compressing activation representations from 16 bits to 8 bits, effectively allowing the context window to be doubled under unchanged hardware constraints. Prospects for further progress include lossless compression to 4 bits, or even the utilization of 1.58-bit representations in future LLMs, which we highlight as a future research direction.

\textbf{LLMs on Edge and Mobile}

Deploying 1.58-bit LLMs promises significant advances for language models operating on edge and mobile hardware. These environments are often severely constrained by memory and compute resources, restricting the scale and effectiveness of deployed LLMs. By lowering both memory and energy requirements, 1.58-bit LLMs make it feasible to bring powerful language capabilities to such platforms, unlocking new applications that were previously unattainable. Additionally, due to the compatibility with standard CPUs—prevalent in edge and mobile ecosystems—\text{BitNet b1.58} can be efficiently run, providing further performance gains and broader applicability for LLM-based solutions on these devices.

\textbf{New Hardware for 1-bit LLMs}

Advancements such as~Groq\footnote{\url{https://groq.com/}} have demonstrated the feasibility and advantages of building tailored hardware solutions (like LPUs) for accelerating LLM workloads. Building upon this direction, we propose and encourage the design and development of novel hardware and systems engineered for 1-bit LLMs, leveraging the unique computational paradigm introduced by BitNet~\cite{bitnet}.